% SPDX-License-Identifier: GPL-3.0-or-later
% Copyright (C) 2026 Alessandro Greco
\chapter{Каналы и протокол соединения}\label{cap:07-canali}
\aggiornatoal{23 сентября 2026 г.}{коммит torspec \texttt{8e1f245}; tor 0.4.9.13; Arti 2.6.0; протокол соединения 4--5}

\begin{obiettivi}
К концу этой главы вы сможете:
\begin{itemize}[nosep]
  \item объяснять, что такое канал и почему Tor строит его поверх TLS;
  \item описывать ключи, образующие идентичность реле, и то, как они связываются между собой;
  \item прослеживать процедуру установления канала ячейка за ячейкой: VERSIONS, CERTS, AUTH\_CHALLENGE, AUTHENTICATE и NETINFO;
  \item читать и проверять сертификат Tor Ed25519;
  \item понимать, почему клиент не аутентифицируется и что это означает для анонимности;
  \item открывать проверенный канал к реле лабораторной сети (лабораторная работа~1).
\end{itemize}
\end{obiettivi}

Цепь, центральная конструкция Tor, представляет собой последовательность переходов. Но каждый переход опирается на более простой элемент: прямое зашифрованное соединение между двумя соседними узлами. Спецификация называет это соединение \textbf{каналом}. В этой главе мы увидим, как устанавливается канал и, прежде всего, как каждая сторона проверяет идентичность того узла, с которым намеревалась соединиться. Формат ячеек, передаваемых по каналу, рассматривается в \cref{cap:08-celle}, к которой мы будем часто обращаться.

% ---------------------------------------------------------------------------
\section{Что такое канал?}\label{sec:canale-def}

Канал --- это прямое зашифрованное соединение между двумя реле либо между клиентом и реле~\spec{tor-spec}{Channels}. Технически это TLS-сессия поверх TCP, переносящая ячейки. Её характеризуют три свойства.

\begin{description}
  \item[Асимметричные роли] Открыть канал может кто угодно (\termine{инициатор}{initiator}), но принять его может только реле (\termine{ответчик}{responder}). Клиент не принимает соединения: у него нет порта, прослушивающего подключения из сети Tor.
  \item[Проверенная идентичность] Ответчик всегда криптографически доказывает владение одной или несколькими идентичностями реле. Инициатор обычно уже знает, какую идентичность должен увидеть: он получает её из консенсуса, из списка доверенных серверов каталогов, встроенного в исходный код, либо, если канал открывается при расширении цепи, из самого сообщения \texttt{EXTEND2}. Если предъявленная идентичность не совпадает с ожидаемой, канал должен быть закрыт.
  \item[Совместное использование] Один канал переносит множество цепей, различаемых по CircID (\cref{sec:circid}). Узлы предпочитают повторно использовать уже существующий канал, а не открывать второй к тому же реле. Канал без цепей закрывается после случайно выбранного периода бездействия, диапазон которого зависит от назначения канала (\file{channelpadding\_get\_channel\_idle\_timeout} в \file{src/core/or/channelpadding.c}): 3--4,5 минуты для клиентских каналов и 45--75 минут для каналов между двумя реле, где центральное значение задаётся параметром консенсуса \texttt{nf\_conntimeout\_relays} (один час). Пока канал используется, \opt{KeepalivePeriod} вместо этого приводит к отправке ячейки заполнения (PADDING) каждые пять минут, чтобы межсетевые экраны не закрывали соединение.
\end{description}

\begin{approfondimento}[Канал или соединение?]
В старой документации и многих научных работах каналы называются \emph{соединениями} (\emph{OR connections}). Спецификация отказалась от этого термина, потому что его можно спутать с TCP-соединениями, которые exit-реле открывают для исходящих потоков. Оба термина сохранились в исходном коде tor: \file{or\_connection\_t} обозначает TLS-соединение, а \file{channel\_t} --- абстракцию, которая его использует.
\end{approfondimento}

% ---------------------------------------------------------------------------
\section{Идентичность реле}\label{sec:canale-chiavi}

Прежде чем разбирать процедуру установления канала, нужно понять, \emph{что именно} должно быть доказано. Реле идентифицируется долговременными ключами. Спецификация использует систематические имена: префикс \file{KP\_} обозначает открытый ключ, а \file{KS\_} --- соответствующий закрытый~\spec{tor-spec}{Relay keys and identities}. Для канала важны следующие ключи:

\begin{description}
  \item[\file{KP\_relayid\_ed}] \textbf{Идентичность Ed25519}, или \enquote{главный ключ идентичности (master identity key)}. Она никогда не истекает и служит одной цели: подписывать ключ подписи. Её можно хранить офлайн.
  \item[\file{KP\_relayid\_rsa}] 1024-битная \textbf{устаревшая идентичность RSA}. На ней по-прежнему основаны 40-значные шестнадцатеричные отпечатки в Tor Metrics, но безопасность протокола больше от неё не зависит.
  \item[\file{KP\_relaysign\_ed}] Среднесрочный \textbf{ключ подписи}, сертифицированный идентичностью. Именно этот ключ реле использует в обычной ежедневной работе онлайн.
  \item[\file{KP\_link\_ed}] Краткосрочный ключ, используемый только тогда, когда реле инициирует канал и должно аутентифицироваться (\cref{sec:canale-authenticate}).
\end{description}

На \cref{fig:chiavi-relay} показано, как сертификаты связывают эти ключи. Идея представляет собой \emph{цепочку доверия}, начинающуюся с идентичности и доходящую до TLS-сертификата конкретного соединения.

\begin{figure}[htbp]
\centering
% SPDX-License-Identifier: GPL-3.0-or-later
% Copyright (C) 2026 Alessandro Greco
%% Русское альтернативное описание EPUB.
\altfig{Цепочка аутентификации ретранслятора Tor: ключ идентичности Ed25519 KP\_relayid\_ed удостоверяет ключ подписи KP\_relaysign\_ed (тип сертификата 4). Ключ подписи удостоверяет дайджест SHA-256 TLS-сертификата отвечающей стороны (тип 5) либо ключ KP\_link\_ed для инициатора, проходящего аутентификацию (тип 6). Устаревший ключ идентичности RSA KP\_relayid\_rsa подписывает перекрёстный сертификат привязки к Ed25519 (тип 7), а также имеет самоподписанный сертификат X.509 (тип 2). Указанные 30 и 2 дня — сроки действия сертификатов в примере, а не неизменная длительность жизни всех ключей; не каждая стрелка означает подпись одной ключевой пары Ed25519 другой.}
% Цепочка сертификатов relay.
\begin{tikzpicture}[font=\sffamily\small,node distance=1.1cm,
  chiave/.style={draw=torviola,fill=torviola!10,rounded corners,minimum width=3.6cm,minimum height=0.9cm,align=center},
  legacy/.style={draw=torgrigio,fill=torgrigio!8,rounded corners,minimum width=3.4cm,minimum height=0.9cm,align=center},
  cert/.style={-{Stealth[length=2.5mm]},thick,torverde!60!black},
  etich/.style={font=\sffamily\scriptsize,fill=white,inner sep=1pt}]
  \node[chiave] (id) {\texttt{KP\_relayid\_ed}\\ключ Ed25519 (долговременный)};
  \node[chiave,below=of id] (sign) {\texttt{KP\_relaysign\_ed}\\ключ подписи (сертиф. 30 дн.)};
  \node[chiave,below left=1.1cm and -1.2cm of sign] (tls) {TLS-сертификат\\SHA-256 (сертиф. 2 дн.)};
  \node[chiave,below right=1.1cm and -1.2cm of sign] (link) {\texttt{KP\_link\_ed}\\аутентификация (сертиф. 2 дн.)};
  \node[legacy,right=2.2cm of id] (rsa) {\texttt{KP\_relayid\_rsa}\\идентичность RSA (устар.)};
  \draw[cert] (id) -- node[etich,right] {4 \texttt{IDENTITY\_V\_SIGNING}} (sign);
  \draw[cert] (sign) -- node[etich,left] {5 \texttt{SIGNING\_V\_TLS\_CERT}} (tls);
  \draw[cert] (sign) -- node[etich,right] {6 \texttt{SIGNING\_V\_LINK\_AUTH}} (link);
  \draw[cert,torgrigio] (rsa) -- node[etich,above] {7} (id);
  \draw[cert,torgrigio] (rsa.north east) to[out=60,in=0,looseness=4] node[etich,above] {2 (самоподписанный)} (rsa.east);
  \node[font=\sffamily\scriptsize,text=torgrigio,below=2pt of tls] {отправляет отвечающая сторона};
  \node[font=\sffamily\scriptsize,text=torgrigio,below=2pt of link] {отправляет инициатор, проходящий аутентификацию};
\end{tikzpicture}

\caption{Цепочка сертификатов, связывающая Ed25519-идентичность реле с TLS-сертификатом соединения. Каждая стрелка обозначает сертификат: ключ у её начала подписывает ключ или отпечаток у конца. Число указывает CertType в ячейке CERTS. Справа показана устаревшая цепочка RSA, связывающая старую RSA-идентичность с Ed25519-идентичностью.}
\label{fig:chiavi-relay}
\end{figure}

В лаборатории эти ключи находятся в файлах каталога \file{keys/} каждого реле: \file{ed25519\_master\_id\_public\_key} и соответствующий секретный ключ идентичности, \file{ed25519\_signing\_cert} и \file{ed25519\_signing\_secret\_key} для ключа подписи, а также \file{secret\_id\_key} для RSA-идентичности. Файл \file{fingerprint-ed25519} в каталоге узла содержит его Ed25519-идентичность в base64.

\begin{storia}[От RSA к Ed25519]
До 2015 года единственной идентичностью реле был его 1024-битный RSA-ключ --- размер, который сегодня считается недостаточным. \proposta{220} ввела Ed25519-идентичности и формат сертификатов, описываемый в этой главе; поддержка появилась в tor 0.2.7.2-alpha. RSA-ключ сохранили ради совместимости и связали с Ed25519-ключом перекрёстным сертификатом (cross-certificate).
\end{storia}

% ---------------------------------------------------------------------------
\section{Процедура установления канала в общих чертах}\label{sec:canale-handshake}

\Cref{fig:handshake-canale} кратко показывает установление канала~\spec{tor-spec}{Negotiating and initializing channels}.

\begin{figure}[htbp]
\centering
% SPDX-License-Identifier: GPL-3.0-or-later
% Copyright (C) 2026 Alessandro Greco
%% Русское альтернативное описание EPUB.
\altfig{Упрощённая схема установления современного канала Tor: инициатор устанавливает TLS, после чего сообщения зашифрованы, а обе стороны обмениваются VERSIONS. Отвечающая сторона последовательно отправляет CERTS, AUTH\_CHALLENGE и NETINFO. Если инициатор — публичный ретранслятор, он обязан аутентифицироваться: отправить CERTS и AUTHENTICATE (шаги в пунктирной рамке). Клиенты и мосты эти два сообщения не отправляют. Затем инициатор в любом случае завершает обмен сообщением NETINFO. Необязательные ячейки VPADDING на схеме опущены; направление стрелок обозначает отправителя.}
% Установление канала.
\begin{tikzpicture}[font=\sffamily\small,>=Stealth]
  \node[client] (i) at (0,0) {Инициатор};
  \node[relay,minimum size=1.3cm] (r) at (7,0) {\scriptsize Ответная сторона};
  \draw[torgrigio] (i.south) -- ++(0,-8.4);
  \draw[torgrigio] (r.south) -- ++(0,-8.4);
  \draw[<->,thick,torgrigio] (0,-1.1) -- node[above]{TLS-рукопожатие} node[below,font=\scriptsize]{начиная отсюда всё зашифровано} (7,-1.1);
  \newcommand{\av}[2]{\draw[cellaf] (0,#1) -- node[above,font=\ttfamily\footnotesize]{#2} (7,#1);}
  \newcommand{\ind}[2]{\draw[risposta] (7,#1) -- node[above,font=\ttfamily\footnotesize]{#2} (0,#1);}
  \av{-2.3}{VERSIONS}
  \ind{-3.0}{VERSIONS}
  \ind{-3.7}{CERTS}
  \ind{-4.4}{AUTH\_CHALLENGE}
  \ind{-5.1}{NETINFO}
  \draw[dashed,torgrigio,rounded corners] (-0.4,-5.55) rectangle (7.4,-7.05);
  \node[font=\scriptsize,text=torgrigio,anchor=east,align=right] at (-0.5,-6.45) {только публичные\\ретрансляторы};
  \av{-6.1}{CERTS}
  \av{-6.8}{AUTHENTICATE}
  \av{-7.8}{NETINFO}
  \node[font=\scriptsize,text=torgrigio,anchor=west] at (7.5,-7.8) {канал открыт};
\end{tikzpicture}

\caption{Процедура установления канала. После установления TLS всё передаётся в зашифрованном виде. Ответчик всегда отправляет CERTS, AUTH\_CHALLENGE и NETINFO именно в таком порядке. Пунктирный блок применяется только тогда, когда инициатор --- публичное реле, которому необходимо подтвердить собственную идентичность; клиент или мост отправляет только NETINFO.}
\label{fig:handshake-canale}
\end{figure}

\begin{enumerate}
  \item Инициатор открывает TLS-соединение.
  \item Обе стороны отправляют ячейку \cella{VERSIONS} и согласуют версию протокола соединения (\cref{sec:celle-formato}).
  \item Ответчик отправляет \cella{CERTS} (сертификаты, подтверждающие идентичность), \cella{AUTH\_CHALLENGE} (вызов, который инициатор подпишет, если захочет аутентифицироваться) и \cella{NETINFO} (время и адреса).
  \item Инициатор проверяет сертификаты. Если он является публичным реле, то отвечает \cella{CERTS} и \cella{AUTHENTICATE}; в любом случае завершает установление канала ячейкой \cella{NETINFO}.
\end{enumerate}

\cella{NETINFO} отмечает конец процедуры установления канала: после её отправки сторона больше не отправляет ячейки установления канала. Так ответчик понимает, решил ли инициатор не аутентифицироваться. Между ними может появляться любое число ячеек \cella{VPADDING}, но никакие другие типы ячеек недопустимы.

% ---------------------------------------------------------------------------
\section{Уровень TLS}\label{sec:canale-tls}

Tor использует TLS необычным образом. У веб-сайта сертификат сервера подписан центром сертификации и содержит доменное имя; браузер проверяет и подпись, и имя. Реле, напротив, предъявляет \textbf{самоподписанный} сертификат без значимого имени, а клиент не проверяет его по обычным правилам TLS. В коде лабораторной работы~1 это выражается двумя строками, которые почти в любом другом контексте были бы серьёзной ошибкой безопасности:
\begin{lstlisting}[language=Python,numbers=none]
ctx.check_hostname = False
ctx.verify_mode = ssl.CERT_NONE
\end{lstlisting}
Здесь это не ошибка, потому что аутентификация выполняется \emph{после}, внутри канала, с помощью ячейки \cella{CERTS}. TLS обеспечивает шифрование и целостность канала и доказывает владение ключом TLS-сертификата; всё остальное устанавливают ячейки.

Спецификация добавляет несколько правил~\spec{tor-spec}{The TLS handshake}:
\begin{itemize}
  \item реализациям следует поддерживать TLS~1.3, а реле --- также TLS~1.2, с современными наборами шифров (AES-GCM, ChaCha20-Poly1305 и обменом ключами X25519 или P-256);
  \item \emph{возобновление TLS-сеанса (session resumption)} не следует разрешать: это ослабляет прямую секретность и уже делало возможными атаки вроде \emph{Triple Handshake} (2013);
  \item сжатие TLS использовать не следует.
\end{itemize}
В лабораторной сети клиент Python с OpenSSL~3.5 согласует с tor TLS~1.3 и набор \file{TLS\_AES\_256\_GCM\_SHA384}.

\begin{storia}[Зачем нужна такая необычная процедура установления канала?]
Самые ранние версии протокола соединения помещали сертификаты Tor непосредственно в TLS-рукопожатие (версия~1), затем скрывали их с помощью повторного согласования TLS (версия~2), сигнализируя выбранный вариант через список предлагаемых наборов шифров. Причиной была цензура: до появления подключаемых транспортов (\cref{cap:16-censura}) Tor пытался выглядеть как обычное HTTPS-соединение между браузером и веб-сервером. В TLS~1.2 сертификаты передавались открыто и легко распознавались. В версии~3 сертификаты переместили внутрь зашифрованного канала, в ячейки \cella{CERTS}. Сегодня сопротивление цензуре возложено на подключаемые транспорты, а старые варианты устарели; спецификация отмечает, что по состоянию на март 2025 года tor всё ещё принимал некоторые из них, но они были запланированы к удалению.
\end{storia}

% ---------------------------------------------------------------------------
\section{CERTS: подтверждение идентичности}\label{sec:canale-certs}

\subsection{Ячейка}

\cella{CERTS} --- ячейка переменной длины (команда 129), тело которой содержит список сертификатов:

\begin{center}
\begin{tabular}{lll}
\toprule
Поле & Байты & Содержимое \\
\midrule
N & 1 & число сертификатов \\
\multicolumn{3}{l}{\emph{повторяется N раз:}} \\
\quad CertType & 1 & тип сертификата \\
\quad CertLen & 2 & длина сертификата \\
\quad Certificate & CertLen & сертификат \\
\bottomrule
\end{tabular}
\end{center}

Ячейка \cella{CERTS} не может содержать два сертификата одного типа. Определённые типы перечислены в \cref{tab:cert-types}~\spec{cert-spec}{Certificate types (CERT\_TYPE field)}.

\begin{table}[htbp]
\centering\small
\begin{tabularx}{\textwidth}{@{}cl l>{\raggedright\arraybackslash}X@{}}
\toprule
Тип & Имя & Формат & Что сертифицирует \\
\midrule
1 & \file{TLS\_LINK\_X509} & X.509 & TLS-ключ, подписанный RSA-идентичностью (устарел) \\
2 & \file{RSA\_ID\_X509} & X.509 & самоподписанная RSA-идентичность \\
3 & \file{LINK\_AUTH\_X509} & X.509 & RSA-ключ аутентификации (устарел) \\
4 & \file{IDENTITY\_V\_SIGNING} & Ed & ключ подписи, подписанный Ed25519-идентичностью \\
5 & \file{SIGNING\_V\_TLS\_CERT} & Ed & отпечаток TLS-сертификата, подписанный ключом подписи \\
6 & \file{SIGNING\_V\_LINK\_AUTH} & Ed & ключ \file{KP\_link\_ed}, подписанный ключом подписи \\
7 & \file{RSA\_ID\_V\_IDENTITY} & RSA$\to$Ed & Ed25519-идентичность, подписанная RSA-идентичностью \\
8--12 & \multicolumn{3}{l}{onion-сервисы (8, 9, 11), перекрёстный сертификат ключа ntor (10), семейства (12)} \\
\bottomrule
\end{tabularx}
\caption{Типы сертификатов. \enquote{Ed} обозначает собственный формат сертификата Tor Ed25519.}
\label{tab:cert-types}
\end{table}

\subsection{Сертификат Tor Ed25519}

Tor не использует X.509 для сертификатов Ed25519. Вместо этого применяется собственный, гораздо более простой формат~\spec{cert-spec}{Tor Ed25519 Certificates}, показанный на \cref{fig:cert-ed}.

\begin{figure}[htbp]
\centering
% SPDX-License-Identifier: GPL-3.0-or-later
% Copyright (C) 2026 Alessandro Greco
%% Русское альтернативное описание EPUB.
\altfig{Структура сертификата Ed25519 сети Tor: VERSION (1 байт), CERT\_TYPE (1), EXPIRATION\_DATE (4: часы с начала эпохи Unix), CERT\_KEY\_TYPE (1), CERTIFIED\_KEY (32: ключ либо дайджест в зависимости от типа), N\_EXTENSIONS (1), затем ноль или более расширений переменной длины и SIGNATURE (64). Каждое расширение содержит ExtLen (2 байта, длина только данных), ExtType (1), ExtFlags (1), ExtData (ExtLen байт). Подпись Ed25519 охватывает все предыдущие поля, но не себя. На схеме имена полей сокращены; начальные смещения: 0, 1, 2, 6, 7, 39 и 40.}
% Сертификат Ed25519 Tor.
\begin{tikzpicture}[campo/.style={draw,minimum height=1.05cm,font=\sffamily\scriptsize,inner sep=1pt,align=center}]
  \node[campo,fill=torviola!35,minimum width=1.1cm] (a) {VER\\(1)};
  \node[campo,fill=torviola!25,minimum width=1.3cm,right=0pt of a] (b) {TYPE\\(1)};
  \node[campo,fill=torviola!15,minimum width=1.8cm,right=0pt of b] (c) {EXPIRATION\\(4)};
  \node[campo,fill=torviola!25,minimum width=1.4cm,right=0pt of c] (d) {KEY\_TYPE\\(1)};
  \node[campo,fill=torverde!15,minimum width=2.4cm,right=0pt of d] (e) {CERTIFIED\_KEY\\(32)};
  \node[campo,fill=torviola!15,minimum width=1.1cm,right=0pt of e] (f) {N\_EXT\\(1)};
  \node[campo,fill=torgrigio!10,minimum width=1.9cm,right=0pt of f] (g) {расширения\\(перем.)};
  \node[campo,fill=avviso!15,minimum width=1.9cm,right=0pt of g] (h) {SIGNATURE\\(64)};
  \foreach \n/\o in {a/0,b/1,c/2,d/6,e/7,f/39,g/40} \node[font=\scriptsize,above] at (\n.north west) {\o};
  \draw[decorate,decoration={brace,mirror,amplitude=5pt}] ([yshift=-2pt]a.south west) -- ([yshift=-2pt]g.south east)
    node[midway,below=6pt,font=\sffamily\small] {байты, покрываемые подписью Ed25519};
  \node[below=1.3cm of g,draw,fill=torgrigio!10,font=\sffamily\scriptsize,align=left] (est)
    {каждое расширение:\\ ExtLen (2) · ExtType (1) · ExtFlags (1) · ExtData (ExtLen)};
  \draw[dotted] (g.south) -- (est.north);
\end{tikzpicture}

\caption{Сертификат Tor Ed25519. Подпись охватывает каждое поле, расположенное перед ней. Срок действия выражается в \emph{часах} с эпохи Unix: беззнаковое четырёхбайтовое поле не переполнится сотни тысяч лет.}
\label{fig:cert-ed}
\end{figure}

\begin{description}
  \item[VERSION] имеет значение 1.
  \item[CERT\_TYPE] имеет то же значение, что CertType в \cref{tab:cert-types}.
  \item[EXPIRATION\_DATE] задаёт срок действия в часах с 1 января 1970 года.
  \item[CERT\_KEY\_TYPE] указывает содержимое следующего поля: 1 означает ключ Ed25519, а 3 --- SHA-256-отпечаток X.509-сертификата. Версии tor до 0.4.5.1-alpha включительно всегда записывали 1, и спецификация требует, чтобы реализации это допускали.
  \item[CERTIFIED\_KEY] содержит 32 сертифицируемых байта.
  \item[Extensions] Каждое расширение содержит длину, тип, флаги и данные. Сегодня определено только расширение типа 4, \emph{signed-with-ed25519-key}, содержащее ключ, подписавший сертификат. Флаг 1 (\file{AFFECTS\_VALIDATION}) определяет, как обрабатывать неизвестное расширение: если он установлен, сертификат необходимо отклонить; иначе расширение игнорируется.
  \item[SIGNATURE] --- подпись Ed25519 всех предыдущих полей.
\end{description}

\subsection{Проверка}\label{sec:canale-verifica}

Ответчик должен отправить ровно один сертификат типа 4 и один типа 5. Инициатор проверяет их следующим образом~\spec{tor-spec}{Authenticating the responder from its CERTS}:
\begin{enumerate}
  \item Сертификат типа 4 содержит в расширении 4 ключ Ed25519: заявленную идентичность \file{KP\_relayid\_ed}. Сертификат должен быть подписан именно этим ключом, то есть быть \emph{самоподписанным}; сертифицированный ключ --- \file{KP\_relaysign\_ed}.
  \item Сертификат типа 5 должен быть подписан \file{KP\_relaysign\_ed}, а его сертифицированным ключом должен быть SHA-256-отпечаток TLS-сертификата, полученного во время TLS-рукопожатия.
  \item Ни один сертификат не должен быть просрочен.
  \item Наконец, полученная идентичность сравнивается с ожидаемой.
\end{enumerate}

Почему эта цепочка что-то доказывает? Шаг 2 связывает ключ подписи с TLS-сертификатом \emph{этого} соединения. TLS-рукопожатие уже доказало, что ответчик владеет закрытым ключом, соответствующим этому сертификату. Шаг 1 связывает ключ подписи с идентичностью. Самозванец мог бы скопировать ячейки \cella{CERTS} настоящего реле, поскольку они не содержат секретов, но не смог бы ими воспользоваться: во время TLS-рукопожатия ему пришлось бы предъявить TLS-сертификат настоящего реле и, следовательно, владеть его закрытым ключом.

\begin{attenzione}[Проверяйте ожидаемую идентичность]
Полностью корректная цепочка доказывает лишь то, что ответчик владеет \emph{какой-то} идентичностью. Если инициатор пропустит шаг~4, любой обладатель любой идентичности сможет выдать себя за требуемое реле. В лабораторной работе~1 параметр \file{expected\_identity} существует именно поэтому, а один из тестов проверяет, что реле с другой идентичностью отклоняется.
\end{attenzione}

\subsection{Устаревшая цепочка RSA}

Инициатор при желании может также проверить RSA-идентичность. Сертификат типа 2 содержит 1024-битный RSA-ключ (самоподписанный, экспонента 65537), а сертификат типа 7 в формате \enquote{перекрёстного сертификата (cross-certificate)} связывает этот ключ с уже проверенной Ed25519-идентичностью. RSA-подпись перекрёстного сертификата вычисляется над SHA-256-отпечатком строки \texttt{"Tor TLS RSA/Ed25519 cross-certificate"}, за которой следуют поля сертификата. Префикс не позволяет повторно использовать ту же подпись в другом контексте.

\subsection{Что реле действительно отправляет}

В лабораторной сети реле \texttt{test003r} отправляет ячейку \cella{CERTS} с пятью сертификатами:
\begin{lstlisting}[numbers=none]
 type 1  584 bytes
 type 2  442 bytes
 type 4  140 bytes  expires 2026-10-24 16:00  key_type=1  ext=[4]
 type 5  104 bytes  expires 2026-09-29 11:00  key_type=3  ext=[]
 type 7  165 bytes
\end{lstlisting}
Помимо Ed25519-цепочки (типы 4 и 5) и RSA-цепочки (типы 2 и 7), реле всё ещё отправляет старый сертификат типа 1 для старых клиентов. Он не просто имеет ту же длину, что TLS-сертификат: совпадает и его отпечаток. Сравнение байтов показывает, что это \emph{точно тот же} сертификат. Его длину не стоит запоминать, и полезно понимать почему. Общее имя (common name) сертификата --- фиктивный домен, который tor строит как \texttt{www.}, затем от 4 до 25 случайных символов и \texttt{.com} (\file{crypto\_random\_hostname} в \file{src/lib/tls/tortls\_openssl.c}). Только это даёт разброс примерно в двадцать байт между реле. Длина меняется и \emph{со временем} у одного реле, потому что tor обновляет TLS-контекст, а вместе с ним сертификат, каждые два часа (\file{MAX\_SSL\_KEY\_LIFETIME\_INTERNAL} в \file{src/core/or/or.h}). Повторное измерение через сорок минут дало тело \cella{CERTS} длиной 1456 байт вместо 1451. Если ваши измерения отличаются, это ожидаемо: совпадать должна \emph{структура}, а сумма длин сертификатов плюс один байт счётчика и три байта заголовка на каждый сертификат должна равняться длине тела ячейки. Две даты истечения показывают разные сроки жизни: ключ подписи сертифицируется на 30 дней, а привязка к TLS-сертификату --- лишь на 2 дня. Это значения tor по умолчанию (параметры \opt{SigningKeyLifetime} и \opt{TestingLinkCertLifetime}, причём второй настраивается только в тестовой сети), поэтому они относятся и к публичной сети. Ключ идентичности остаётся постоянным, тогда как представляющие его онлайн-ключи часто обновляются, ограничивая ущерб от компрометации реле.

% ---------------------------------------------------------------------------
\section{AUTH\_CHALLENGE и AUTHENTICATE}\label{sec:canale-authenticate}

\subsection{Кто аутентифицируется?}

Аутентификация инициатора необязательна, а правило, определяющее, кто должен её выполнять, относится к важнейшим архитектурным решениям протокола~\spec{tor-spec}{Channels}:
\begin{itemize}
  \item \textbf{Публичное реле}, открывающее канал к другому реле, \emph{должно, как правило}, аутентифицироваться: тогда ответчик знает, что разговаривает с сетевым реле, а не клиентом.
  \item \textbf{Клиент} или \textbf{мост} \emph{не должен, как правило}, аутентифицироваться.
\end{itemize}
Причина второго правила --- анонимность. У клиента нет идентичности, и её создание было бы вредно: guard-реле могло бы узнавать того же клиента при повторном подключении с другого адреса. Мост, в свою очередь, не должен раскрывать своим пользователям, что является сетевым реле. Компромисс заключается в том, что реле может сделать вывод, что оно является первым переходом цепи, если канал не аутентифицирован (\cref{cap:08-celle}, \cref{sec:relay}). Это одна из причин, по которым реле отказываются выступать exit-реле для одношаговых цепей.

\subsection{Вызов}

Ответчик в любом случае отправляет ячейку \cella{AUTH\_CHALLENGE}. Она содержит 32 случайных байта и список допустимых методов аутентификации:

\begin{center}
\begin{tabular}{lll}
\toprule
Поле & Байты & Содержимое \\
\midrule
Challenge & 32 & случайное значение, отдельное для каждого канала \\
N\_Methods & 2 & число методов \\
Methods & $2 \times$ N\_Methods & список методов \\
\bottomrule
\end{tabular}
\end{center}

Определены методы 1 (\texttt{RSA-SHA256-TLSSecret}, устарел), 2 (никогда не был реализован) и 3 (\texttt{Ed25519-SHA256-RFC5705}). Ячейка лабораторного реле имеет тело длиной 36 байт: $32 + 2 + 2$, что соответствует одному методу, номер 3.

\subsection{Ответ}

Реле, желающее аутентифицироваться, отвечает собственной ячейкой \cella{CERTS} (используя типы 4 и 6 вместо 4 и 5, поскольку инициатор не обязательно предъявлял TLS-сертификат), после чего отправляет \cella{AUTHENTICATE}. Для метода 3 тело \cella{AUTHENTICATE} представляет собой 352-байтовый документ, подписанный ключом \file{KP\_link\_ed}:

\begin{center}\small
\begin{tabularx}{\textwidth}{@{}l c >{\raggedright\arraybackslash}X@{}}
\toprule
Поле & Байты & Содержимое \\
\midrule
TYPE & 8 & строка \texttt{AUTH0003} \\
CID, SID & 32 + 32 & SHA-256-отпечатки RSA-идентичностей инициатора и ответчика \\
CID\_ED, SID\_ED & 32 + 32 & Ed25519-идентичности инициатора и ответчика \\
SLOG & 32 & SHA-256 всех байтов, \emph{полученных} по каналу до AUTH\_CHALLENGE включительно \\
CLOG & 32 & SHA-256 всех байтов, \emph{отправленных} по каналу до AUTHENTICATE \\
SCERT & 32 & SHA-256 TLS-сертификата ответчика \\
TLSSECRETS & 32 & значение, выведенное из секретов TLS-сессии (\emph{keying material exporter}) \\
RAND & 24 & случайные байты \\
SIG & 64 & подпись Ed25519 всех предыдущих полей \\
\bottomrule
\end{tabularx}
\end{center}

Каждое поле связывает подпись с определённым аспектом \emph{этого} соединения: участвующими идентичностями (CID, SID и т.\,д.), предыдущим диалогом (SLOG и CLOG, включающими случайный вызов), TLS-сертификатом (SCERT) и самой TLS-сессией (TLSSECRETS). Такую подпись нельзя повторно использовать ни в другом соединении, ни посреднику, пересылающему сообщения между двумя разными TLS-сессиями (атака \emph{man-in-the-middle}), поскольку TLSSECRETS и SCERT на двух сторонах будут различаться.

\begin{approfondimento}[Ошибка, сохранённая в спецификации]
TLSSECRETS вычисляется через TLS \emph{exporter} (RFC~5705 и, для TLS~1.3, RFC~8446, §\,7.5) с меткой \texttt{"EXPORTER FOR TOR TLS CLIENT BINDING AUTH0003"} и контекстным значением. Предполагалось использовать в качестве контекста Ed25519-идентичность инициатора, однако реализация по ошибке использовала CID, то есть отпечаток RSA-идентичности. Поскольку это не ослабляет безопасность, спецификация решила документировать фактическое поведение, а не менять его. Это хороший пример того, как для протокола с тысячами уже развёрнутых реализаций поведение реального кода может стать стандартом.
\end{approfondimento}

% ---------------------------------------------------------------------------
\section{NETINFO: время и адреса}\label{sec:canale-netinfo}

Последняя ячейка процедуры установления канала, отправляемая обеими сторонами, --- \cella{NETINFO}. Это ячейка фиксированной длины (команда 8) со следующим телом:

\begin{center}
\begin{tabularx}{\textwidth}{@{}ll>{\raggedright\arraybackslash}X@{}}
\toprule
Поле & Байты & Содержимое \\
\midrule
TIME & 4 & секунды с эпохи Unix, big-endian \\
OTHERADDR & 2 + ALEN & адрес другой стороны глазами отправителя \\
NMYADDR & 1 & число собственных адресов отправителя \\
MYADDR & NMYADDR $\times$ (2 + ALEN) & собственные адреса отправителя \\
\bottomrule
\end{tabularx}
\end{center}

Каждый адрес кодируется типом (один байт: \texttt{0x04} для IPv4, \texttt{0x06} для IPv6), длиной (4 или 16) и значением. Эти поля решают две задачи:
\begin{itemize}
  \item \textbf{Время} позволяет обнаруживать неверные часы: tor пишет предупреждение, если время, полученное от реле, отличается от собственного более чем на час (константа \file{NETINFO\_NOTICE\_SKEW} в \file{src/core/or/channeltls.c}).
  \item \textbf{Адреса} позволяют реле узнать публичный адрес, под которым оно видно, а инициатору --- проверить, достиг ли он реле по одному из его \emph{канонических} адресов, то есть адресов, опубликованных в дескрипторе.
\end{itemize}

\begin{attenzione}[Клиент не раскрывает свои часы]
Реле отправляют собственное время и адреса. Клиент, напротив, должен отправлять \texttt{00 00 00 00} как время и не должен отправлять ни одного собственного адреса. Причина --- \emph{идентификация по отпечатку}: часы компьютера почти всегда имеют небольшой характерный сдвиг. Если клиент сообщал бы время каждому guard-реле, этот сдвиг мог бы помочь узнавать клиента с течением времени. Это иллюстрирует общее правило, к которому мы ещё не раз вернёмся: клиент Tor не должен раскрывать ничего, что не является строго необходимым.
\end{attenzione}

Беззнаковое четырёхбайтовое поле TIME переполнится в 2106 году; \proposta{338} рассматривает его расширение до восьми байт. Получатель ячейки \cella{NETINFO} не должен слепо доверять значениям: другая сторона может солгать о времени и адресах.

Вот ячейка \cella{NETINFO}, полученная от реле \texttt{test003r} и декодированная лабораторной работой~1:
\begin{lstlisting}[numbers=none]
Netinfo(timestamp=1790511791, peer_address='127.0.0.1',
        my_addresses=['127.0.0.1'])
\end{lstlisting}
Время соответствует 27 сентября 2026 года, 12:23:11 UTC. Оба адреса являются адресами обратной петли (loopback), потому что вся сеть работает на одном компьютере. Наш ответ, построенный \file{client\_netinfo\_cell}, имеет тело \texttt{00000000 04 04 7f000001 00}: нулевое время, адрес реле и ни одного собственного адреса.

% ---------------------------------------------------------------------------
\section{После установления канала}\label{sec:canale-dopo}

Теперь канал \emph{открыт} и может переносить цепи. Простой эксперимент обнаруживает интересную деталь. Если после открытия канала слушать его, не создавая цепей, реле не отправляет \emph{ничего}: в лаборатории за пятнадцать секунд не приходит ни одной ячейки.

Можно было бы ожидать ячейки \cella{PADDING}: начиная с версии~5 протокола соединения клиенты и guard-реле обмениваются ячейками заполнения на бездействующих каналах через случайные интервалы от 1,5 до 9,5 секунды, чтобы сетевой наблюдатель, например маршрутизатор, записывающий статистику NetFlow, не мог отличить периоды активности от тишины~\spec{padding-spec}{Connection-level padding}. Исходный код tor объясняет, почему здесь этого не происходит. Функция, решающая, добавлять ли заполнение (\file{channelpadding\_decide\_to\_pad\_channel} в \file{src/core/or/channelpadding.c}), активирует его только на каналах, \emph{уже использованных} для полных цепей или пользовательского трафика. Со стороны реле канал считается использованным так только тогда, когда через него проходит многошаговая цепь: первая ячейка \cella{RELAY\_EARLY}, расширяющая цепь ко второму переходу, либо \cella{RELAY} в цепи, продолжающейся за пределы реле (\file{src/core/or/relay.c}). Новый канал остаётся безмолвным, как и канал, несущий только одношаговые цепи, что мы увидим в \cref{cap:09-circuiti}. Заполнение появится в \cref{cap:10-cifratura-relay}, когда наш клиент расширит цепь.

Этот небольшой эпизод даёт методологический урок: \emph{отсутствие наблюдения не является доказательством}. Прежде чем заключать, что заполнение не работает, нужно проверить, выполнены ли условия, при которых он вообще должен появляться.

% ---------------------------------------------------------------------------
\section{Лабораторная работа: открытие проверенного канала}\label{sec:lab-canale}

\begin{laboratorio}[1]{Открыть канал к реле и проверить его идентичность}
\textbf{Цель.} Написать \file{labs/ru/код/torlab/channel.py}, который открывает канал как клиент Tor: TLS, \cella{VERSIONS}, проверка \cella{CERTS}, чтение \cella{AUTH\_CHALLENGE} и \cella{NETINFO} и отправка собственной \cella{NETINFO}.\\
\textbf{Предварительные условия.} Лабораторная работа~2 (\file{cells.py}, \cref{cap:08-celle}): обработка ячеек лежит в основе канала, поэтому логично сначала написать её, хотя объясняется она в следующей главе. Здесь нисходящее объяснение (сначала канал, затем ячейки) и восходящая реализация (сначала ячейки, затем канал) идут в противоположных направлениях, и это нормально. Также необходима работающая лабораторная сеть (\cref{cap:19-rete-privata}) и библиотека \texttt{cryptography} для Ed25519.\\
\textbf{Оценочное время.} 2--3 часа.
\end{laboratorio}

\paragraph{Шаг 1: декодировать сертификат Ed25519.} Следуя \cref{fig:cert-ed}, читаем заголовок фиксированной длины, затем расширения и сохраняем байты, охваченные подписью:

\begin{lstlisting}[language=Python]
def decode_ed_certificate(data):
    if len(data) < 40 + 64 or data[0] != 1:
        raise ChannelError("некорректный сертификат")
    _, kind, expiration, key_type = struct.unpack(">BBIB", data[:7])
    key = data[7:39]
    extension_count, pos, extensions = data[39], 40, {}
    for _ in range(extension_count):
        length, extension_type, flag = struct.unpack(">HBB", data[pos:pos+4])
        body = data[pos+4:pos+4+length]
        if extension_type != 4 and flag & 1:     # AFFECTS_VALIDATION
            raise ChannelError("неизвестное расширение влияет на проверку")
        extensions[extension_type] = (flag, body)
        pos += 4 + length
    return EdCertificate(kind, expiration, key_type, key,
                         extensions, data[:pos], data[pos:])
\end{lstlisting}

\paragraph{Шаг 2: проверить цепочку.} Эта функция применяет правила \cref{sec:canale-verifica} и возвращает доказанную идентичность:

\begin{lstlisting}[language=Python]
def verify_responder_identity(certificates, cert_tls_der):
    id_v_sign = decode_ed_certificate(certificates[4])
    sign_v_tls = decode_ed_certificate(certificates[5])
    identity = id_v_sign.signing_key               # extension 4
    id_v_sign.verify_signature(identity)             # self-signed
    sign_v_tls.verify_signature(id_v_sign.certified_key)
    if sign_v_tls.certified_key != sha256(cert_tls_der).digest():
        raise ChannelError("TLS-сертификат не совпадает")
    for c in (id_v_sign, sign_v_tls):
        if c.expired():
            raise ChannelError("сертификат просрочен")
    return identity
\end{lstlisting}

TLS-сертификат в DER-формате получаем из сокета вызовом \file{sock.getpeercert(binary\_form=True)}. Python возвращает его даже несмотря на то, что в данном случае мы его ещё не проверяли.

\paragraph{Шаг 3: установление канала.} Функция \file{open\_channel} отправляет \cella{VERSIONS}, предлагая версии 4 и 5 (как Arti, мы пропускаем 3). Затем она читает ячейки до появления \cella{NETINFO}, обрабатывая каждый тип: игнорирует \cella{VPADDING} и все \cella{VERSIONS} после первой, а любую другую неожиданную ячейку отклоняет. Наконец, проверяет сертификаты, сравнивает идентичность с ожидаемой и отправляет собственную \cella{NETINFO}. Полный код находится в репозитории; его центральный цикл:

\begin{lstlisting}[language=Python]
while netinfo is None:
    cell, buffer = _read(sock, buffer, version)
    if cell.command == CellCommand.VERSIONS:
        if version:
            continue
        version = negotiate_version(SUPPORTED_VERSIONS, cell.body)
    elif cell.command == CellCommand.CERTS:
        certificates = decode_certs(cell.body)
    elif cell.command == CellCommand.AUTH_CHALLENGE:
        ...                                    # допустимые методы
    elif cell.command == CellCommand.NETINFO:
        netinfo = decode_netinfo(cell.body)
    elif cell.command != CellCommand.VPADDING:
        raise ChannelError("неожиданная ячейка во время handshake")
\end{lstlisting}

\paragraph{Шаг 4: тесты без сети.} В \file{labs/ru/код/torlab/test/test\_channel.py} генерируются два ключа Ed25519, строится корректная цепочка сертификатов и проверяется, что код её принимает. Затем цепочка нарушается всеми способами, описанными в спецификации: изменённая подпись, несертифицированный ключ подписи, другой TLS-сертификат, истёкший сертификат, повторяющийся тип; для каждого дефекта проверяется обнаружение. Тест только корректного случая ничего не доказывает: проверяющая функция, принимающая всё подряд, тоже его пройдёт.

\paragraph{Шаг 5: запуск в лабораторной сети.} При работающей сети:
\begin{lstlisting}[language=Python]
from torlab.channel import open_channel, identity_from_file

expected = identity_from_file("labs/ru/лаборатория/net/nodes/003r/fingerprint-ed25519")
c = open_channel("127.0.0.1", 5103, expected_identity=expected)
print(c.version, c.ed_identity.hex(), c.netinfo, c.auth_methods)
print(c.registry)
\end{lstlisting}

Сокращённый вывод:
\begin{lstlisting}[numbers=none]
5 785d3199f21d20ad...  Netinfo(time=1790511793, ...)  [3]
['TLS TLSv1.3 TLS_AES_256_GCM_SHA384', '<- VERSIONS (6 байт)',
 '<- CERTS (1451 байт)', '<- AUTH_CHALLENGE (36 байт)',
 '<- NETINFO (509 байт)', '-> NETINFO: канал открыт']
\end{lstlisting}
Если вместо этого подключиться к порту 5104 (реле \texttt{test004r}), ожидая идентичность другого реле, получается корректная ошибка: \texttt{реле имеет идентичность, отличную от ожидаемой}.

% ---------------------------------------------------------------------------
\section*{Итоги}
\begin{itemize}
  \item Канал --- это TLS-сессия между двумя соседними узлами. Каналы принимают только реле, а один канал переносит множество цепей.
  \item Идентичность реле --- его Ed25519-ключ \file{KP\_relayid\_ed}, сертифицирующий ключ подписи \file{KP\_relaysign\_ed}; RSA-идентичность сохраняется ради совместимости.
  \item TLS-сертификат реле самоподписан. Аутентификация происходит внутри канала через ячейку \cella{CERTS}, связывающую идентичность, ключ подписи и отпечаток TLS-сертификата.
  \item Инициатор обязан проверить подписи, сроки действия, привязку к TLS и \textbf{ожидаемую идентичность}.
  \item Реле, инициирующие канал, аутентифицируются с помощью \cella{AUTHENTICATE}, подписи, привязанной к TLS-сессии и диалогу. Клиенты и мосты не аутентифицируются, чтобы не создавать узнаваемых идентичностей.
  \item \cella{NETINFO} завершает процедуру установления канала. Клиент отправляет нулевую временную метку и ни одного собственного адреса.
  \item Заполнение на уровне канала требует, чтобы канал уже переносил \emph{полную} цепь или пользовательский трафик. Новый канал или канал только с одношаговыми цепями остаётся безмолвным. После прохождения этого порога заполнение запускается также до первого прикладного трафика (\texttt{nf\_pad\_before\_usage}, по умолчанию 1).
\end{itemize}

% ---------------------------------------------------------------------------
\section*{Упражнения}

\begin{esercizio}[1]\label{es:07-ruoli}
Почему клиенту Tor не нужен прослушивающий порт, доступный из сети? В каких случаях реле инициирует каналы?
\end{esercizio}

\begin{esercizio}[1]\label{es:07-scadenza}
Сертификат Ed25519 имеет EXPIRATION\_DATE, равную \texttt{0x000772F7}. Когда истекает его срок? Почему спецификация использует часы, а не секунды?
\end{esercizio}

\begin{esercizio}[1]\label{es:07-authchallenge}
Тело ячейки \cella{AUTH\_CHALLENGE} имеет длину 38 байт. Сколько методов она предлагает? Корректна ли 36-байтовая ячейка с N\_Methods, равным 2?
\end{esercizio}

\begin{esercizio}[2]\label{es:07-impostore}
Противник записывает настоящую ячейку \cella{CERTS} реле и без изменений воспроизводит её на собственном сервере, предъявляя собственный TLS-сертификат. Какой шаг проверки завершится ошибкой? Что изменится, если противник также предъявит TLS-сертификат настоящего реле?
\end{esercizio}

\begin{esercizio}[2]\label{es:07-netinfo}
Запишите в шестнадцатеричном виде тело ячейки \cella{NETINFO}, отправленной реле с адресом \texttt{203.0.113.5} в 12:00:00 UTC 1 января 2026 года клиенту, который подключился с адреса \texttt{198.51.100.7}.
\end{esercizio}

\begin{esercizio}[2]\label{es:07-client-auth}
Предположим, клиенты Tor аутентифицируются перед guard-реле ключом, который генерируется один раз и хранится на диске. Опишите два конкретных способа, которыми это повредило бы анонимности пользователя.
\end{esercizio}

\begin{esercizio}[2]\label{es:07-rsa}
Расширьте \file{labs/ru/код/torlab/channel.py}, добавив необязательную проверку RSA-идентичности (типы 2 и 7) по \cref{sec:canale-certs}. Сравните SHA-1-отпечаток полученного RSA-ключа с файлом \file{fingerprint} реле. \emph{Подсказка:} отпечаток --- SHA-1 DER-кодировки открытого ключа.
\end{esercizio}

\begin{esercizio}[3]\label{es:07-padding}
При работающей лабораторной сети откройте канал с помощью \file{labs/ru/код/torlab/channel.py} и держите его открытым, пока клиент \texttt{test008c} загружает страницу через сеть. Получите ли вы ячейки \cella{PADDING} по \emph{своему} каналу? Сформулируйте гипотезу, проверьте её и объясните результат с учётом \file{channelpadding\_decide\_to\_pad\_channel}.
\end{esercizio}

\begin{esercizio}[3]\label{es:07-mitm}
Объясните, почему поле TLSSECRETS в \cella{AUTHENTICATE} не позволяет посреднику ретранслировать аутентификацию реле $A$ реле $B$ через другую TLS-сессию, даже если посредник точно пересылает каждую ячейку. Что произошло бы без TLSSECRETS?
\end{esercizio}

% ---------------------------------------------------------------------------
\section*{Дополнительная литература}
\begin{itemize}
  \item \titolo{Tor Protocol Specification}: главы \enquote{Channels}, \enquote{Negotiating and initializing channels} и \enquote{Relay keys}, а также \titolo{Tor certificate formats} (cert-spec)~\cite{torspec}.
  \item Предложение 220, \titolo{Migrate server identity keys to Ed25519}: причины перехода от RSA.
  \item \titolo{Tor channel заполнение} в спецификации заполнение: модель противника на основе NetFlow и выбор интервалов.
  \item RFC~5705 и RFC~8446 (§\,7.5): функции TLS \titolo{keying material exporter}.
\end{itemize}
